\subsection{Batatia et al.\ (2024): MACE-MP-0 foundation model}
\label{sec:appendix-example-batatia2024mp0}

\paragraph{Q1 -- Bibliographic identification.}
Batatia, Benner, Chiang, et al.\ (Cs\'anyi group) release MACE-MP-0,
an equivariant message-passing foundation model trained on the MPtrj
dataset and demonstrated zero-shot across solids, liquids, gases,
surfaces, MOFs, and a small protein. arXiv:2401.00096.
\lstinputlisting[language=turtle]{artifacts/kg/listings/batatia2024mp0/q01-bibliographic.ttl}

\paragraph{Q2 -- Material system.}
The MLIP covers MPtrj's ${\sim}150{,}000$ Materials Project structures
across 89 elements, with 90\% of unit cells under 70 atoms. We model
this as a single $\MaterialSystemC$; the paper's out-of-distribution
demonstrations (water, organics, MOFs, protein) are documented in
Q12.
\lstinputlisting[language=turtle]{artifacts/kg/listings/batatia2024mp0/q02-material.ttl}

\paragraph{Q3 -- Reference calculation method.}
Reference data come from DFT calculations.
\lstinputlisting[language=turtle]{artifacts/kg/listings/batatia2024mp0/q03-reference-method-type.ttl}

\paragraph{Q4 -- Reference settings.}
DFT uses VASP with PAW pseudopotentials and PBE GGA, with Hubbard $U$
applied on transition-metal oxides containing Co/Cr/Fe/Mn/Mo/Ni/V/W
combined with O or F (Materials Project default $U$-values). No
dispersion correction is used in training.
\lstinputlisting[language=turtle]{artifacts/kg/listings/batatia2024mp0/q04-reference-settings.ttl}

\paragraph{Q5 -- Training dataset.}
${\sim}1.5$\,M static and structural-relaxation configurations across
${\sim}150{,}000$ MP compounds, covering energies, forces, and
stresses. Provenance is \emph{published} (re-using the CHGNet-curated
MPtrj split).
\lstinputlisting[language=turtle]{artifacts/kg/listings/batatia2024mp0/q05-dataset.ttl}

\paragraph{Q6 -- Sampling strategies.}
Two strategies were combined: MP relaxation-trajectory sampling
(reusing the CHGNet/Deng et al.\ 2023 curated set) and periodic-table
chemical-space coverage (89 elements; binary, ternary, and
higher-order chemistries).
\lstinputlisting[language=turtle]{artifacts/kg/listings/batatia2024mp0/q06-sampling.ttl}

\paragraph{Q7 -- MLIP method, components, implementation.}
MACE-MP-0 is an equivariant message-passing graph network (e3nn) with
two layers of higher-order equivariant messages built from sums of
two-body permutation-invariant polynomials in a spherical basis;
correlation order 3 (4-body messages per layer); a tensor-decomposition
parameterisation; and a ZBL repulsive pair potential at close range
with Agnesi distance transform. The MACE-MP-0b3 release is the L=1
medium-sized variant in $128\!\times\!0e + 128\!\times\!1o$ irreps.
Loss is a weighted Huber on energy/forces/stresses with weights
$(1, 10, 10)$ and a force Huber-$\delta$ that decays with force
magnitude. Optimisation is AMSGrad with EMA learning-rate schedule.
Implementation: mace-mp (PyTorch + e3nn).
\lstinputlisting[language=turtle]{artifacts/kg/listings/batatia2024mp0/q07-method.ttl}

\paragraph{Q8 -- Hyperparameter settings.}
Reported settings: cutoff $r_c = 6.0$\,\AA, 2 message-passing layers,
$l_{\max}=3$, correlation order 3, 128 channels, 10 radial Bessel
basis functions, $3{\times}64$ SiLU readout MLP, initial learning rate
$10^{-3}$, 100 epochs.
\lstinputlisting[language=turtle]{artifacts/kg/listings/batatia2024mp0/q08-hyperparameter-settings.ttl}

\paragraph{Q9 -- MLIP run and trained model.}
A single training run on MPtrj produces the MACE-MP-0b3 medium
foundation model. Training compute is reported on the run.
\lstinputlisting[language=turtle]{artifacts/kg/listings/batatia2024mp0/q09-run-and-model.ttl}

\paragraph{Q10 -- Benchmark results.}
Test-set MAEs on the MPtrj held-out split: 18\,meV/atom on energy and
39\,meV/\AA{} on force components. Many additional zero-shot
out-of-distribution evaluations (organic-molecule MD, MOF flexibility,
protein side-chain stability) are reported but not encoded as
separate $\BenchmarkResultC$ instances.
\lstinputlisting[language=turtle]{artifacts/kg/listings/batatia2024mp0/q10-benchmark-results.ttl}

\paragraph{Q11 -- Computational resources.}
Reported: 2{,}600 GPU-hours of training on 40--80 NVIDIA H100 GPUs
across 10--20 nodes; inference scaling on a single A100 80GB GPU
(${\sim}$ns/day for a 1000-atom system; weak scaling to 32{,}000 atoms
on 64 GPUs at 0.1\,ns/day). Wall-clock training duration, peak
memory, and explicit per-atom inference time are not reported.
\lstinputlisting[language=turtle]{artifacts/kg/listings/batatia2024mp0/q11-resources.ttl}

\paragraph{Q12 -- Gaps.}
The most consequential gaps are: \emph{out-of-distribution use}
($\dlConcept{TrainedModel}$ is applied zero-shot to organics, MOFs, a
protein, and gas-phase chemistry far beyond the MPtrj training set,
but the ontology has no concept for "deployment domain" distinct from
the training $\MaterialSystemC$); explicit DFT settings that vary by
MP task (energy cutoff, k-mesh) are not surfaced; and the foundation
model's three released sizes (small/medium/large) plus the L=1/L=2
variants are not encoded as separate $\TrainedModelC$ instances. The
training duration in wall-clock time is also not reported, only GPU
hours.
